\subsection{MODULES OF FINITE LENGTH}

Recall (I, § 4, no. 4, Definition 7) that an A-module M is called simple if it is not reduced to 0 and it contains no submodule distinct from M and \{0\}. An A-module M is said to be of finite length if it has a Jordan-Hölder series \((\mathbf{M}_t)_{0 \leqslant t \leqslant n}\) and the number \(n\) of quotients of this series (which does not depend on the Jordan-Hölder series of M considered) is then called the length of M (I, § 4, no. 7, Definition 11); we shall denote it by long(M) or long\_A(M). An A-module which is reduced to 0 is of length 0; if M is a non-zero A-module of finite length then long(M) \textgreater{} 0.

PROPOSITION 16. Let M be an A-module and N a submodule of M; for M to be of finite length, it is necessary and sufficient that N and M/N be so, and then

\[
\operatorname{long} (\mathbf {N}) + \operatorname{long} (\mathbf {M} / \mathbf {N}) = \operatorname{long} (\mathbf {M}).\tag{34}
\]

The proof has been given in I, § 4, no. 7, Proposition 10.

COROLLARY 1. Let M be an A-module of finite length; for a submodule N of M to be equal to M, it is necessary and sufficient that long(N) = long(M).

COROLLARY 2. Let \(u: \mathbf{M} \to \mathbf{N}\) be an A-module homomorphism. If \(\mathbf{M}\) or \(\mathbf{N}\) is of finite length, so is \(\operatorname{Im}(u)\) . If \(\mathbf{M}\) is of finite length, so is \(\operatorname{Ker}(u)\) and

\[
\operatorname{long} (\operatorname{Im} (u)) + \operatorname{long} (\operatorname{Ker} (u)) = \operatorname{long} (\mathbf {M}).\tag{35}
\]

If \(\mathbf{N}\) is of finite length, so is \(\operatorname{Coker}(u)\) and

\[
\operatorname{long} (\operatorname{Im} (u)) + \operatorname{long} (\operatorname{Coker} (u)) = \operatorname{long} (\mathbf {N}).\tag{36}
\]

COROLLARY 3. Let \((\mathbf{M}_i)_{0\leqslant i\leqslant n}\) be a finite family of A-modules of finite length. If there exists an exact sequence of linear mappings

\[
0 \longrightarrow \mathrm{M} _ {0} \stackrel {u _ {0}} {\longrightarrow} \mathrm{M} _ {1} \stackrel {u _ {1}} {\longrightarrow} \mathrm{M} _ {2} \longrightarrow \dots \longrightarrow \mathrm{M} _ {n - 1} \stackrel {u _ {n - 1}} {\longrightarrow} \mathrm{M} _ {n} \longrightarrow 0\tag{37}
\]

then

\[
\sum_ {k = 0} ^ {n} (- 1) ^ {k} \operatorname{long} (\mathbf {M} _ {k}) = 0.\tag{38}
\]

The corollary is obvious for \(n = 1\) and is just Proposition 16 for \(n = 2\) ; we argue by induction on \(n\) . If \(\mathbf{M}_{n-1}' = \operatorname{Im}(u_{n-2})\) , then, by the induction hypothesis,

\[
\sum_ {k = 0} ^ {n - 2} (- 1) ^ {k} \operatorname{long} (\mathbf {M} _ {k}) + (- 1) ^ {n - 1} \operatorname{long} (\mathbf {M} _ {n - 1} ^ {\prime}) = 0.
\]

On the other hand, the exact sequence \(0 \to \mathbf{M}_{n-1}' \to \mathbf{M}_{n-1} \to \mathbf{M}_n \to 0\) gives

\[
\operatorname{long} \left(\mathbf {M} _ {n - 1} ^ {\prime}\right) + \operatorname{long} \left(\mathbf {M} _ {n}\right) = \operatorname{long} \left(\mathbf {M} _ {n - 1}\right),
\]

whence relation (38).

COROLLARY 4. Let M and N be two submodules of finite length of an A-module E; then M + N is of finite length and

\[
\operatorname{long} (\mathbf {M} + \mathbf {N}) + \operatorname{long} (\mathbf {M} \cap \mathbf {N}) = \operatorname{long} (\mathbf {M}) + \operatorname{long} (\mathbf {N}).\tag{39}
\]

It suffices to apply Corollary 3 to the exact sequence (29) (no. 7).

\[
0 \rightarrow \mathbf {M} \cap \mathbf {N} \rightarrow \mathbf {M} \oplus \mathbf {N} \rightarrow \mathbf {M} + \mathbf {N} \rightarrow 0
\]

using the fact that \(\operatorname{long}(\mathbf{M} \oplus \mathbf{N}) = \operatorname{long}(\mathbf{M}) + \operatorname{long}(\mathbf{N})\) by (34).

COROLLARY 5. Let M be an A-module the sum of a family (N \(_{t}\) ) of submodules of finite length. Then M is of finite length and

\[
\operatorname{long} (\mathbf {M}) \leqslant \sum_ {i} \operatorname{long} \left(\mathrm{N} _ {i}\right).\tag{40}
\]

Moreover, for the two sides of (40) to be equal, it is necessary and sufficient that M be the direct sum of the \(N_{t}\) .

It has been seen (no. 7, formula (28)) that there is a canonical surjective linear mapping \(h \colon \bigoplus_{i} N_{i} \to M\) ; hence the corollary follows from (35).

COROLLARY 6. Let M and N be two submodules of an A-module E such that E/M and E/N are modules of finite length; then E/(M ∩ N) is of finite length and

\[
\text { long } (\mathbf {E} / (\mathbf {M} \cap \mathbf {N})) + \text { long } (\mathbf {E} / (\mathbf {M} + \mathbf {N})) = \text { long } (\mathbf {E} / \mathbf {M}) + \text { long } (\mathbf {E} / \mathbf {N}). \tag {41}
\]

It suffices to apply Corollary 3 to the exact sequence (30)

\[
0 \rightarrow \mathrm{E} / (\mathrm{M} \cap \mathrm{N}) \rightarrow (\mathrm{E} / \mathrm{M}) \oplus (\mathrm{E} / \mathrm{N}) \rightarrow \mathrm{E} / (\mathrm{M} + \mathrm{N}) \rightarrow 0
\]

using the fact that

\[
\operatorname{long} ((\mathbf {E} / \mathbf {M}) \oplus (\mathbf {E} / \mathbf {N})) = \operatorname{long} (\mathbf {E} / \mathbf {M}) + \operatorname{long} (\mathbf {E} / \mathbf {N}).
\]

COROLLARY 7. Let \((\mathbf{M}_i)\) be a finite family of submodules of an A-module \(\mathbf{E}\) such that the \(\mathbf{E} / \mathbf{M}_i\) are modules of finite length. Then \(\mathbf{E} / (\bigcap_{i} \mathbf{M}_i)\) is of finite length and

\[
\operatorname{long} \left(\mathrm{E} / \left(\bigcap_ {i} \mathrm{M} _ {i}\right)\right) \leqslant \sum_ {i} \operatorname{long} (\mathrm{E} / \mathrm{M} _ {i}).\tag{42}
\]

It has been seen (formula (27)) that there is a canonical injective linear mapping \(\mathbf{E} / \left(\bigcap_{i}\mathbf{M}_{i}\right)\to \prod_{i}(\mathbf{E} / \mathbf{M}_{i})\)

Remark. With the exception of no. 7, Proposition 9, all the results of nos. 2 to 10 are valid for arbitrary commutative groups with operators, submodules (resp. quotient modules) being replaced in the statements by stable subgroups (resp. quotient groups by stable subgroups); we also make the convention of calling homomorphisms of groups with operators ``linear mappings''. The corollaries to no. 7, Proposition 9 are also valid for commutative groups with operators: this is obvious for Corollaries 4 and 5, as also for Corollary 2, since \(\alpha\left(\sum_{i\in I}x_i\right)=\sum_{i\in I}\alpha x_i\) for every operator \(\alpha\) , and Corollary 3 follows from it. As for Corollary 1, it suffices to note that if N is a stable subgroup of F containing \(u(S)\) , \(\bar{u}^{-1}(N)\) is a stable subgroup of E containing S, hence \(\bar{u}^{-1}(N)\) contains M and therefore \(u(M)\subset N\) .

